\documentclass[11pt,a4paper]{article}
\usepackage[utf8]{inputenc}
\usepackage[T1]{fontenc}
\usepackage{lmodern}
\usepackage{geometry}
\usepackage{amsmath,amssymb,bm}
\usepackage{hyperref}
\geometry{margin=2.2cm}

\title{Carbon Cycle Numerical Model: Implementation Report}
\author{}
\date{\today}

\begin{document}
\maketitle

\section{Scope}
This report documents what is implemented in \texttt{carbone.py} and how it is implemented:
\begin{itemize}
    \item an 8-state carbon-cycle ODE model,
    \item a forward-Euler time integration scheme,
    \item post-processing and plotting with per-curve min--max normalization,
    \item automatic PDF naming with encoded initial conditions, horizon, and time step.
\end{itemize}

\section{State, Parameters, and Forcing}
The state vector is ordered as
\[
\bm{x}(t)=
\begin{bmatrix}
A & R & D & F & P & S & O & V
\end{bmatrix}^{\!\top},
\]
where:
\begin{itemize}
    \item $A$: atmosphere carbon,
    \item $R$: carbonate rock carbon,
    \item $D$: deep ocean carbon,
    \item $F$: fossil-fuel carbon stock,
    \item $P$: plant carbon,
    \item $S$: soil carbon,
    \item $O$: surface ocean carbon,
    \item $V$: vegetated land area (\%).
\end{itemize}

Useful constants (from the script) are:
\[
\alpha_A=\frac{280}{A_0},\quad
k_P=\frac{55}{P_0},\quad
k_S=\frac{55}{S_0},\quad
k_{sed}=\frac{0.1}{D_0},\quad
k_{down}=\frac{90.1}{O_0},\quad
k_{up}=\frac{90}{D_0}.
\]
With deforestation parameter $\mathrm{Def}$:
\[
\beta=\frac{\mathrm{Def}}{2},
\qquad
\delta=\left(\frac{\mathrm{Def}}{P_0}\cdot 0.2\right)\,100,
\qquad
\nu=0.1 \text{ (volcanic flux)}.
\]

The fossil-combustion forcing is piecewise linear in time from tabulated data:
\[
\mathcal{F}_{comb}(t)=\text{linear interpolation of }(t_i, f_i),
\]
and applied with stock saturation logic:
\[
u_F(t,F)=
\begin{cases}
\mathcal{F}_{comb}(t), & F>0,\\
0, & F\le 0.
\end{cases}
\]

\section{ODE System in Matrix Form (Linear + Nonlinear Split)}
The implemented system can be written as
\[
\dot{\bm{x}} = A_{\mathrm{lin}}\,\bm{x} + \bm{b} + \bm{g}(\bm{x},t).
\]

\subsection*{Linear part}
\[
A_{\mathrm{lin}}=
\begin{bmatrix}
-\mathrm{Kao}\,\alpha_A & 0 & 0 & 0 & k_P & k_S & 0 & 0\\
0 & 0 & k_{sed} & 0 & 0 & 0 & 0 & 0\\
0 & 0 & -(k_{sed}+k_{up}) & 0 & 0 & 0 & k_{down} & 0\\
0 & 0 & 0 & 0 & 0 & 0 & 0 & 0\\
0 & 0 & 0 & 0 & -2k_P & 0 & 0 & 0\\
0 & 0 & 0 & 0 & k_P & -k_S & 0 & 0\\
\mathrm{Kao}\,\alpha_A & 0 & k_{up} & 0 & 0 & 0 & -k_{down} & 0\\
0 & 0 & 0 & 0 & 0 & 0 & 0 & 0
\end{bmatrix},
\]
\[
\bm{b}=
\begin{bmatrix}
\beta+\nu\\
-\nu\\
0\\
0\\
-2\beta\\
\beta\\
0\\
-\delta
\end{bmatrix}.
\]

\subsection*{Nonlinear / time-dependent part}
\[
\bm{g}(\bm{x},t)=
\begin{bmatrix}
\nu_F(t,F)+\mathrm{Kao}\,pCO_2^{oc}(A,O)-\Phi(A,V)\\
0\\
0\\
-\nu_F(t,F)\\
\Phi(A,V)\\
0\\
-\mathrm{Kao}\,pCO_2^{oc}(A,O)\\
0
\end{bmatrix}.
\]

Where:
\begin{itemize}
    \item $\Phi(A,V)$ is photosynthesis,
    \[
    \Phi(A,V)=110\cdot CO2Effect(\alpha_A A)\cdot\frac{V}{100}\cdot TempEffect\!\left(GlobalTemp(\alpha_A A)\right),
    \]
    \item $pCO_2^{oc}(A,O)$ is the ocean partial pressure chain
    \[
    A\to GlobalTemp\to WaterTemp\to (K_{CO_2},K_{carb}),\quad O\to SurfCConc\to HCO_3^-\to CO_3^{2-}\to pCO_2^{oc},
    \]
    which is nonlinear (rational terms and square root in carbonate chemistry).
\end{itemize}

Hence, the model is \textbf{affine-linear plus nonlinear closure/forcing}. The main nonlinearities are $\Phi(A,V)$ and $pCO_2^{oc}(A,O)$, plus the piecewise term $\nu_F(t,F)$.

\section{Numerical Time Integration}
The solver is explicit Euler:
\[
\bm{x}_{n+1}=\bm{x}_n+\Delta t\,\bm{f}(\bm{x}_n,t_n),
\]
with time grid generated by
\[
t_n=t_0+n\Delta t,\quad n=0,\dots,N,\quad t_N\approx t_f.
\]
In code:
\begin{itemize}
    \item \texttt{derivative(x,t)} computes $\bm{f}(\bm{x},t)$,
    \item \texttt{step(x,t,dt)} applies one Euler step,
    \item \texttt{run\_simulation(...)} iterates over all steps and stores trajectories.
\end{itemize}

\section{Plotting and Output Encoding}
Implemented plotting features:
\begin{itemize}
    \item 3 active subplots in a 2\,$\times$\,2 canvas (top-left, top-right, bottom-left),
    \item per-curve min--max normalization
    \[
    y_{norm}=\frac{y-y_{\min}}{y_{\max}-y_{\min}},
    \]
    with constant series mapped to $0.5$,
    \item legends include each curve's physical range $[y_{\min},y_{\max}]$ and use text color matching each curve,
    \item PDF filename encodes initial conditions, simulated years, and $\Delta t$ in scientific notation-safe tags.
\end{itemize}

\section{Reproducibility Notes}
The report corresponds to the current \texttt{carbone.py} structure and parameterization (including current deforestation value in the script). If parameters or flux functions change, update the matrix split accordingly.

\end{document}
